\section{Decisión de agente}

Se eligió \texttt{ToolCallingAgent} porque la tarea consiste en pocas consultas conocidas. Sus llamadas tienen argumentos definidos y quedan registradas por separado, lo que facilita validarlas y revisar la traza~\cite{smolagentsTour,smolagentsSec}. La traza es el registro de las acciones, las respuestas de las herramientas y la salida final.

\texttt{CodeAgent} permitiría combinar consultas y cálculos mediante código generado. Sería útil para secuencias menos previsibles, pero exigiría ejecutar ese código en un entorno aislado y con recursos limitados. Las restricciones del intérprete local no proporcionan ese aislamiento~\cite{smolagentsSec}. Se dejó esta opción fuera de la evaluación. La Figura~\ref{fig:acciones} compara ambas formas de expresar la tarea y distingue el pseudocódigo de las llamadas realmente ejecutadas.

\begin{figure}[htbp]
\centering
\begin{minipage}[t]{0.47\textwidth}
\footnotesize\textbf{Pseudocódigo, no ejecutado}
\begin{verbatim}
filas = [resultados('demo', 'es', s)
         for s in ('A', 'B')]
mejor = min(filas, key=lambda f: f['wer'])
final_answer(json.dumps({...}))
\end{verbatim}
\end{minipage}\hfill
\begin{minipage}[t]{0.5\textwidth}
\footnotesize\textbf{Traza real con llamadas estructuradas}
\begin{verbatim}
resultados(corpus=demo, idioma=es, sistema=A)
resultados(corpus=demo, idioma=es, sistema=B)
final_result_Comparacion(ganador=B, ...)
\end{verbatim}
\end{minipage}
\caption{La misma consulta como código y como llamadas estructuradas. La izquierda es ilustrativa. La derecha resume la traza \texttt{pydantic\_mcp\_caso1\_r1}, cuyas llamadas solicitó el modelo.}
\label{fig:acciones}
\end{figure}
